\newpage

\section{What Is Data?}
\label{sec:what-is-data}

A simple example is probably the best place to start. Look at the
cross-section of a tree trunk and you find facts there without anyone needing
to explain what a tree even is. The number of growth rings gives a reasonable
estimate of the tree's age, since each ring corresponds to one year of growth
laid down during the annual cycle of dormancy and activity. The texture and
color of the bark hint at the species, and warped grain, discoloration, or
insect tunnels point to its health. Look closer still, and the density of the
wood, the direction of the grain, and the moisture content explain the piece's
weight and volume. None of these facts were created the moment someone noticed
them. They were already there in the tree, available to anyone who knew how to
read them. That correspondence, between a physical state and a fact about the
world you can read off from it, is about the simplest illustration there is of
what this paper means by \emph{data}.

\subsection{Etymology and Lexical History}

The word \emph{data} is the plural of the Latin \emph{datum}, the past
participle of \emph{dare}, ``to give.'' A datum is, quite literally, a
``given'': a piece of information taken as a starting point for reasoning
rather than derived from something else. In everyday speech the plural form
gets treated as a singular collective noun all the time (``the data is
clear''), but in technical writing the plural is preferred: data \emph{are}
observations, measurements, or recorded facts, not one undivided thing.\\
The Latin origin carries real structural weight that survives into technical
usage. A datum is something received from the world, not something invented
by the observer who happens to record it. Standard references preserve this
sense, but they never quite formalize it.
\newpage
\subsection{Existing Definitions and Their Shared Structure}

Different sources define the term in different, only partly overlapping ways.
Rob, Morris, and Coronel describe data as \emph{raw facts}, facts that have not
yet been processed enough for their meaning to become
apparent~\autocite{rob2013}. Ratzan instead defines the companion term:
\emph{information} is a coherent collection of data, organized in a particular
way and meaningful as a result~\autocite{ratzan2004}. The two definitions
agree on more than they disagree about. In both, data come before meaning
rather than supplying it directly.\\
General dictionaries describe roughly the same shape from a different angle.
Webster's \textit{New International Dictionary} treats data as something given
or accepted: facts or principles granted, the basis from which an inference or
a line of reasoning proceeds, or the material an idealized system is built
from~\autocite{websters1934}. The \textit{Oxford English Dictionary} is more
compact about it: data are facts or things known, used as the basis for
inference or calculation~\autocite{oed2023}. The UNESCO definition adds a
structural requirement that the others leave out:

\begin{quote}
Facts, concepts, or instructions expressed in a standardized way, suitable for
being communicated, interpreted, or processed, whether by people or by
automated systems.~\autocite{unesco2008}
\end{quote}

``Expressed in a standardized way'' is really the key phrase here. It gestures
at, without actually specifying, the requirement that data have some
\emph{structure} that makes indexing, comparison, and transmission possible in
the first place. The informal definitions all agree that data precede meaning,
but none of them says what structural properties make data addressable to
begin with.
\newpage
\subsection{The Structural Implication}

Put these definitions together and they point at something none of them states
outright. Data have no shape that belongs to them independently of some
particular point of view. A fact only counts as data once it sits inside a
framework, whether of meaning, purpose, or use, relative to which it is
relevant, organized, coherent, and useful. A growth ring means nothing to a
system that has no concept of seasons; a voltage means nothing to a circuit
built to ignore it. Data presuppose some prior agreement, however informal,
about what is even worth attending to.\\
This structural implication is exactly what the formal definition in
Section~\ref{sec:formal-definition} pins down. Before getting there, it is
worth seeing the same implication play out in a setting further removed from
bytes and circuits, and then, closer to engineering, in one setting where the
stakes are more familiar.

\subsection{A Second Illustration: The Same Text, Different Data}
\label{sub:second-illustration}

The bytes example below is deliberately narrow. A single passage of ordinary
prose makes the same point at a larger scale, and shows that
framework-dependence is not an artifact of binary encodings, it is a property
of data as such. Consider the opening of Marx's account of the 1851 coup:

\begin{quote}
Upon the overthrow of the first Napoleon came the restoration of the Bourbon
throne (Louis XVIII, succeeded by Charles X). In July, 1830, an uprising of
the upper tier of the bourgeoisie, or capitalist class, the aristocracy of
finance, overthrew the Bourbon throne, or landed aristocracy, and set up the
throne of Orleans, a younger branch of the house of Bourbon, with Louis
Philippe as king. From the month in which this revolution occurred, Louis
Philippe's monarchy is called the ``July Monarchy.'' In February, 1848, a
revolt of a lower tier of the capitalist class, the industrial bourgeoisie,
against the aristocracy of finance, in turn dethroned Louis Philippe. This
affair, also named from the month in which it took place, is the ``February
Revolution.''~\autocite{marx1852}
\end{quote}
\newpage Nothing about this passage changes as it is read. It is one fixed string of
characters. What changes is what different readers, bringing different
frameworks, are able to extract from it as data:

\begin{description}
  \item[A reader with no background in the subject] can recover only the
        coarsest structure: that the passage is in English, that it names
        people called Napoleon and Louis, and that a couple of four-digit
        numbers, 1830 and 1848, appear in it.
  \item[A reader competent in the language but not the subject] adds
        linguistic structure on top of that: sentence patterns, a register of
        vocabulary that reads as historical and political, without yet being
        able to say what any of it refers to.
  \item[A historian] recovers named events: the fall of the Bourbon
        restoration, the July Monarchy, the February Revolution, and the
        succession of political transitions in France between roughly 1815
        and 1848.
  \item[A political economist] recovers a different structure from the exact
        same sentences: two distinct fractions of the bourgeoisie, financial
        and industrial, in conflict over political power, with the landed
        aristocracy displaced earlier in the sequence.
  \item[A political philosopher] reads the passage as data about the state,
        class, and the mechanics of revolutionary transformation, largely
        independent of which specific individuals are named.
  \item[A literary scholar] attends to none of the above directly, and
        instead extracts data about diction, sentence construction, and how
        the passage represents historical causation through its syntax.
\end{description}
Each reader is looking at the identical string of characters. None of them is
wrong. What each one is able to treat as data is fixed entirely by the
framework they bring to the text, exactly as the structural implication above
predicts: no fact becomes data until it is read inside a framework relative
to which it is relevant and organized.\\
The passage also supports a distinction worth making explicit. Some of what a
historian recovers is stated directly, for instance that Louis Philippe's
monarchy took its name from the month of July. Some of it is not stated but
follows from analysis, such as the relation between the earlier fall of the
Bourbon throne and the later fall of the July Monarchy, a relation the
sentences never announce but that a competent reader constructs from them.
Both count as data under the definition adopted here; the second kind is
simply data obtained by inference over the first, rather than read off the
surface directly.\newpage
One detail in the passage raises a question this paper does not resolve here.
Nothing about the concept of July connects intrinsically to the monarchy that
took its name from it, and nothing about February connects intrinsically to
the revolution named after it. A name with no inherent tie to the thing it
names was adopted, and every subsequent reader accepts that name without
friction, using it to pick out the event unambiguously. The bytes example
above shows that the same signal can carry different data depending on
framework; this example shows the converse problem, that a token can come to
stand reliably for something with which it shares no intrinsic connection at
all. That is a question about \emph{representation} rather than about data as
defined here, and it is a question this paper returns to once the formal
apparatus of Section~\ref{sec:formal-definition} is available to state it
precisely.
\newpage
\subsection{The Context-Dependence of Data}

The same physical signal can be different data depending on the framework
applied to it. Take five bytes with hexadecimal values \texttt{0x41 0x52 0x4C
0x49 0x5A}. Under one set of assumptions this is the ASCII string
\texttt{"ARLIZ"}; under another it is five unsigned 8-bit integers, 65, 82, 76,
73, and 90; under a third it is the first five bytes of some arbitrary binary
file format; under a fourth, a fragment of a network packet payload. The bytes
themselves are identical every time. What changes is not the underlying
signal but the interpretation imposed on top of it. Meaning is not contained in
the bytes. It arises from the agreement between whatever system produced them
and whatever system reads them.\\
This example also shows the practical cost of ignoring context-dependence: if
the framework is never stated, different components will end up applying
different frameworks to the same bytes, and they will produce different,
incompatible results. The defects described in
Section~\ref{sec:introduction} are exactly this phenomenon, just at scale.

\subsection{The Perishability of Data}

A common assumption early in a technical career is that data are fixed and
objective, that once you record a fact, the fact just sits there. This is only
partly true, and the part that isn't true matters quite a lot. Walliman points
out that data are not only imprecise but \emph{transient} and
\emph{perishable}~\autocite{walliman2011}. A poll measuring voting intentions
ahead of an election will not produce the same result as an otherwise
identical poll run at a different time, even with the same sampling method
and, in principle, the same respondents.\\
``Perishable'' does not mean the data physically vanish. A measurement
recorded yesterday is still sitting there to be read today. What perishes is
its \emph{currency}, its claim to describe the present state of whatever it
measured. Data perish in a second sense too. Secondhand reports and biased
accounts get presented as settled fact all the time, because no method of
recording information is perfectly reliable, and some distortion creeps in at
almost every stage of transmission. This is not some rare edge case confined to
opinion polling. It is a basic property of how facts move from the world into
a recorded form. The engineering implication is that reliable data require
explicit attention to precision, integrity, provenance, and timestamp, and
these properties follow naturally once you adopt the formal definition in
Section~\ref{sec:formal-definition}.